%%%%%%%%%%%%%%%%%%%%%%%%%%%%%
%% Chapter 1: The Introduction %%
%%%%%%%%%%%%%%%%%%%%%%%%%%%%%

% DRAFT by Claude — Dion to revise; argument framing is mine to approve

\section{Introduction}
\label{chap:intro}

M dwarfs are the most numerous stars in the Galaxy and the most common planet hosts, which places them at the centre of the search for habitable worlds \citep{Bonfils_2013, Dressing_2015, Gaidos_2016, Henry_2018}. Their low luminosities bring the habitable zone close in \citep{Kopparapu_2013}, where planets are easier to detect \citep{Charbonneau_2007} but also more exposed to stellar activity \citep{Shields_2016}. M dwarfs remain magnetically active for billions of years, and the resulting high-energy radiation and particle flux drive star--planet interactions \citep{Vidotto_2018} that shape the atmospheres and habitability of the planets they host \citep{Hebrard_2016, Dumusque_2021, Bellotti_2022}. A growing body of work has measured and characterised these small-scale magnetic fields \citep{Saar_1985, Johns_1996, shulyak+2014, Kochukhov_2021, reiners+2022, cristofari_2023a, cristofari_2023b, Donati_2025}.

Most of the magnetic flux on an M dwarf resides in small-scale, mixed-polarity structures that circular polarimetry cannot recover, because opposite polarities cancel in the net Stokes V signal \citep{donati_landstreet+2009, Kochukhov_2021}. These photospheric fields nonetheless leave a measurable imprint on the unpolarised Stokes I spectrum: when a magnetic field is present, energy levels split into sublevels, resulting in broadening of magnetically sensitive lines \citep{Landi_2004, reiners_basri_2007, Reiners+2012, shulyak+2014}. Because these fields reshape the spectrum, neglecting them biases the atmospheric parameters derived from it, while poorly constrained atmospheric parameters in turn bias the inferred field. Measuring the surface magnetic fields of M dwarfs is therefore a prerequisite for both stellar and exoplanetary characterisation.

Both quantities are accessible from the same data. Previous studies \citep{Rajpurohit_2018, Passegger_2019, Marfil_2021, Sarmento_2021} have used high-resolution near-infrared spectrometers, such as CRIRES+ \citep{Crires}, CARMENES \citep{Carmenes}, or SPIRou \citep{spirou}, to constrain the atmospheric properties of tens to hundreds of stars, and the same observations make it possible to measure the field itself by modelling the shape of the lines it broadens \citep{kochukhov+2020, Petit_2021}. Recent works, like \citet{cristofari_2023a, cristofari_2023b, reiners+2022, Hahlin_2023, Petit_2021}, have used the Zeeman broadening as a quantitative tool by fitting M-dwarf spectra against grids of magnetic synthetic spectra. Because the Zeeman effect scales with the effective Land\'e factor, lines of different magnetic sensitivity respond differently to the same field, which helps disentangle magnetic broadening from rotation and macroturbulence. The detection of Zeeman signatures nonetheless remains difficult for weakly magnetised stars, whose lines may be dominated by rotational or macroturbulent broadening. 

These analyses fix the grid of magnetic components in advance and tailor their line selection by hand or through fixed temperature cuts. This thesis instead applies a single homogeneous optical analysis to the whole sample, combining a per-star line-list filter with a data-driven selection of the number of magnetic components, so that the magnetic model of each star is set by its own spectrum rather than fixed beforehand.

The running title of this thesis promises a self-consistent retrieval. We use \emph{self-consistent} to mean that the atmospheric parameters and the magnetic filling factors are retrieved jointly from the same spectra in a single posterior, rather than fixing one while fitting the other.
% TODO(Dion): confirm this is the intended meaning of self-consistent

The thesis is organised as follows. Sect.~\ref{sec:observations} describes the spectroscopic data and their reduction. Sect.~\ref{chap:zeeturbo} introduces the \texttt{ZeeTurbo} magnetic synthetic-spectrum grids. Sect.~\ref{sec:methods} sets out the \texttt{asap} retrieval, the per-star line-list filter, and the BIC-based magnetic-component selection. Sect.~\ref{sec:results} reports the retrieved atmospheric parameters and magnetic fields and compares them with published values. Sect.~\ref{chap:discussion} discusses the results, and Sect.~\ref{chap:conclusions} concludes.
